\subsection{Mémoire profonde et compatibilité des cylindres}
\label{memprof:seccion}

Le prolongement compare des états qui conservent simultanément restes,
quotients et frontières. Dans la coordinatisation linéaire à six coordonnées du registre,
nous prenons le représentant entier de \(L_0\) dont les entrées ont été indiquées :
\[
 A_H=\begin{pmatrix}
 2&2&2&1&2&1\\2&1&2&2&1&1\\1&0&0&1&0&2\\
 0&0&1&1&1&0\\2&2&2&0&2&1\\2&1&2&1&0&0
 \end{pmatrix},\qquad\det A_H=1.
\]
Le choix de ce représentant fait partie de la coordinatisation : la réduction
modulo trois conserve \(L_0\), tandis que le représentant entier fixe
l'opération sur les quotients profonds. Avec des vecteurs-lignes, on obtient
\begin{equation}
 \begin{aligned}
 \mathsf{Div}_H:\mathbb Z^6&\longrightarrow
       \{0,1,2\}^6\times\mathbb Z^6,\\
 x_H&\longmapsto(u_H,x_H^+),\\
 y_H&=x_HA_H,\qquad u_H=[y_H]_3,\qquad
 x_H^+=(y_H-u_H)/3 .
 \end{aligned}
 \label{memprof:division}
\end{equation}
La paire \((u_H,x_H^+)\) conserve \(y_H=u_H+3x_H^+\). Puisque \(A_H\) est unimodulaire, son inverse est entière et
\begin{equation}
 x_H=(u_H+3x_H^+)A_H^{-1}.
 \label{memprof:inversa}
\end{equation}
Cette identité récupère la coordonnée profonde à partir de son résidu et de son quotient. Les autres champs de l'état conservent leurs propres registres. L'orientation ternaire est lue ensuite au moyen de \(0\mapsto0,\ 1\mapsto+1,\ 2\mapsto-1\), en conservant le résidu \(u_H\in\{0,1,2\}^6\) et le quotient de \eqref{memprof:division}.

La coordinatisation cylindrique relie un préfixe ternaire de longueur \(L\),
d'indice entier \(N\), à \(k_{\mathrm{dec}}\) triades décimales,
d'indice \(D\). Les cylindres semi-ouverts correspondants sont
\[
 I_3(N,L)=\left[\frac N{3^L},\frac{N+1}{3^L}\right),\qquad
 I_{1000}(D,k_{\mathrm{dec}})
 =\left[\frac D{1000^{k_{\mathrm{dec}}}},
         \frac{D+1}{1000^{k_{\mathrm{dec}}}}\right).
\]
Nous définissons les entiers de frontière
\[
 a_{\mathrm{cyl}}=N1000^{k_{\mathrm{dec}}}-D3^L,\qquad
 b_{\mathrm{cyl}}=(N+1)1000^{k_{\mathrm{dec}}}-D3^L.
\]
Pour un bloc de six résidus, nous fixons son codage entier
\[
 \nu(u_H)=\sum_{j=1}^6(u_H)_j3^{6-j}
 \in\{0,\ldots,728\}.
\]

\begin{lemma}[Mise à jour entière des cylindres compatibles]
\label{memprof:cilindros}
Les cylindres précédents ont une intersection non vide si et seulement si
\(a_{\mathrm{cyl}}<3^L\) et \(b_{\mathrm{cyl}}>0\).
En ajoutant un bloc ternaire d'indice
\(\nu\in\{0,\ldots,728\}\), on a \(N'=729N+\nu\) et
\(L'=L+6\). Pour \(\Delta k_{\mathrm{dec}}\in\{0,1\}\),
les frontières sont mises à jour par
\[
 \begin{array}{ll}
 \Delta k_{\mathrm{dec}}=0:
   &a_{\mathrm{cyl}}'=729a_{\mathrm{cyl}}+\nu1000^{k_{\mathrm{dec}}},\\[2pt]
 \Delta k_{\mathrm{dec}}=1:
   &D'=1000D+\delta,\\
   &a_{\mathrm{cyl}}'=729000a_{\mathrm{cyl}}
      +\nu1000^{k_{\mathrm{dec}}+1}-\delta\,729\,3^L ,
 \end{array}
\]
où \(\delta\in\{0,\ldots,999\}\) est la triade du prolongement
compatible. Dans les deux cas,
\[
 b_{\mathrm{cyl}}'=a_{\mathrm{cyl}}'
          +1000^{k_{\mathrm{dec}}+\Delta k_{\mathrm{dec}}}.
\]
\end{lemma}
\begin{proof}
Deux intervalles semi-ouverts s'intersectent précisément lorsque chaque extrémité gauche est inférieure à l'extrémité droite de l'autre. Multiplier ces deux inégalités par \(3^L1000^{k_{\mathrm{dec}}}\) donne \(a_{\mathrm{cyl}}<3^L\) et \(b_{\mathrm{cyl}}>0\). La substitution de \(N'\), \(L'\) et \(D'\) dans
\[
 a_{\mathrm{cyl}}'
 =N'1000^{k_{\mathrm{dec}}+\Delta k_{\mathrm{dec}}}-D'3^{L'}
\]
produit les deux récurrences. Soustraire les définitions de
\(b_{\mathrm{cyl}}'\) et \(a_{\mathrm{cyl}}'\) donne la dernière identité.
\end{proof}

Ces opérations agissent sur des préfixes entiers et leurs restrictions
compatibles. La lecture archimédienne apparaît lorsqu'on évalue le système de
cylindres résultant. Le domaine de prolongement incorpore en outre
l'orientation, la retenue, la signature conjointe et les données de frontière de l'état
\(\widehat X\) de \eqref{eq:actualizacion} ; la comparaison des intervalles
est une composante de cette compatibilité conjointe.
